\subsubsection{アーキテクチャの用途}
アーキテクチャの第一の用途は、システムに関わる人々が共有理解を持つことである。共有理解がなければ、設計者、開発者、テスト担当、運用担当が異なる前提で作業し、結果として整合しないシステムになりやすい。

第二の用途は、設計と構築の指針を与えることである。アーキテクチャは、どの構成要素が存在するか、どの構成要素がどの責任を持つか、どのように通信するか、どの品質特性をどの方法で満たすかを示す。これにより、実装前に代替案を比較し、リスクを見つけられる。

第三の用途は、既存システムの理解である。保守、拡張、移行、改善を行うとき、既存システムのアーキテクチャを理解することが重要である。これは、逆方向の設計理解、すなわちリバースアーキテクティングに近い活動である。

また、アーキテクチャは組織構造とも関係する。コンウェイの法則は、システムの構造が、それを設計する組織のコミュニケーション構造を反映しやすいことを指摘する。組織とアーキテクチャの対応は、うまく使えば分担と責任を明確にする力になるが、悪く働くと不要な分断や複雑な依存関係を生む。

\begin{statementbox}{実務での見方}
アーキテクチャは、計画、見積り、分担、再利用、製品系列化にも使われる。単一システムの内部構造だけでなく、複数製品に共通する土台を作るためにも重要である。
\end{statementbox}
